%% -*- mode: TeX -*-
%% -*- coding: utf-8 -*-
\documentclass[UTF8,zihao=-4]{ctexart} %%for macOS set: fontset=macnew
%% ==模式切换 =====
\usepackage[final]{mnsfc} %%draft is default, final is for the final version
%% 加载 amsrefs 包，可以在这里控制选项
\usepackage[alphabetic,abbrev,lite,msc-links]{amsrefs}  %% 如果需要启用引文页码, 添加 backrefs 选项
%% 如需更好地支持 arXiv 参考文献（eprint、超链接、手工标签等），
%% 推荐配合 AMSRefs arXiv Support 使用：
%%   https://github.com/vanabel/AMSRefs-arXiv-Support
%% 并在加载 amsrefs 之后添加（可直接复制到自己的申请书项目中）：
%% \BibSpec{misc}{%
%%   +{} {\PrintAuthors} {author}
%%   +{,} { \textit} {title}
%%   +{,} { } {date}
%%   +{,} { } {eprint}
%%   +{,} { } {note}
%%   +{.} {} {transition}
%% }

\mnsfcProjectTitle{（填写项目名称）}{（Project Title）}
\mnsfcApplicationCode{申请代码1}{申请代码2}  %% 申请代码
\mnsfcKeywordsZH{关键词1；关键词2；关键词3}
\mnsfcKeywordsEN{keywords1; keywords2; keywords3}
\mnsfcResearchField{主要研究领域}  %% 主要研究领域（可选，默认为空）
%% 选择研究属性（二选一）：
%% \mnsfcResearchTypeFree  %% 自由探索类基础研究
%% \mnsfcResearchTypeTarget %% 目标导向类基础研究
%% 或者使用自定义文本：
%% \mnsfcResearchType{自由探索类基础研究}
%% \mnsfcResearchType{目标导向类基础研究}
\mnsfcResearchTypeFree  %% 示例：选择自由探索类基础研究
%% 如需在英文摘要结束后将“报告正文”起始页设为第 7 页，可在导言区启用：
%% \mnsfcBodyStartPage{7}

 \begin{document}
   \mnsfcTitlePage
   \begin{mnsfcAbstractZH}
     这里是\emph{最终可提交版本}的中文摘要内容。
   \end{mnsfcAbstractZH}

   \begin{mnsfcAbstractEN}
     This is the English abstract for the NSFC proposal.
   \end{mnsfcAbstractEN}

  \mnsfcNote{正文开始...}
  \zhlipsum[1-2]
  %% 使用示例：subsubsection、paragraph、enumerate、itemize
   \subsubsection{三级标题示例}
   这是三级标题下的正文内容。三级标题使用编号格式 (1), (2), (3)...

   \paragraph{段落标题示例}
   这是段落标题下的正文内容。段落标题不带编号，仅作为强调使用。

   \subsubsection{列表环境示例}

   \textbf{有序列表（enumerate）示例：}
   \begin{enumerate}
     \item 第一层第一项：编号格式为 1), 2), 3)...
     \item 第一层第二项
     \begin{enumerate}
       \item 第二层第一项：编号格式为 a), b), c)...
       \item 第二层第二项
       \begin{enumerate}
         \item 第三层第一项：编号格式为 i), ii), iii)...
         \item 第三层第二项
       \end{enumerate}
     \end{enumerate}
     \item 第一层第三项
   \end{enumerate}

   \textbf{无序列表（itemize）示例：}
   \begin{itemize}
     \item 第一层第一项：使用实心圆点 •
     \item 第一层第二项
     \begin{itemize}
       \item 第二层第一项：使用空心圆点 ○
       \item 第二层第二项
       \begin{itemize}
         \item 第三层第一项：使用短横线 -
         \item 第三层第二项
       \end{itemize}
     \end{itemize}
     \item 第一层第三项
   \end{itemize}
   \subsubsection{参考文献}
   标准的参考文献数据可以从\href{https://mathscinet.ams.org/mrlookup}{美国数学会(MathSciNet)}获取。
   \cite{EinsteinBergmann1938AM}以及
   \cites{EinsteinBergmann1938AM,Hamilton1982JDG,Osserman1990AMM}
   \section[（为什么要开展此项研究，研究的科学技术价值如何）]{立项依据：}
   \zhlipsum[3-4]
   \begin{equation}\label{eq:NL}
     \int_a^b f'(x)\,dx=f(b)-f(a).
   \end{equation}
  从 \eqref{eq:NL}知,\ldots
  %%\bibliographystyle{plain}%
  \nocite{*}%
   \mnsfcReferences{mnsfc-refs}

   \section[（提纲不做限制，请按照研究工作的自身逻辑撰写。应提炼出特色与创新点、年度研究计划）]{研究内容：}
   \mnsfcNote{建议用"科学问题 → 国际现状 → 关键缺口 → 本项目目标"结构}
   \zhlipsum[5-6]
   \section{研究基础：}

   \subsection[（与本项目相关的研究工作积累和已取得的研究工作成绩，研究风险的应对措施等）；]{研究基础与可行性分析}
   \zhlipsum[7-8]

   \subsection[（包括已具备的实验条件，尚缺少的实验条件和拟解决的途径，包括利用国家实验室、全国重点实验室和部门重点实验室等研究基地的计划与落实情况）；]{工作条件}
   \zhlipsum[9-15]
   \subsection[（申请人和主要参与者正在承担的与本项目相关的科研项目情况，包括国家自然科学基金的项目和国家其他科技计划项目，要注明项目的资助机构、项目类别、批准号、项目名称、获资助金额、起止年月、与本项目的关系及负责的内容等）；]{正在承担的与本项目相关的科研项目情况}
  \mnsfcExcludeLine{无。}
  \subsection[（对申请人负责的前一个已资助期满的科学基金项目（项目名称及批准号）完成情况、后续研究进展及与本申请项目的关系加以详细说明。另附该项目的研究工作总结摘要（限500字）和相关成果详细目录）。]{完成国家自然科学基金项目情况}
  \mnsfcExcludeLine{无。}
  \section{其他需要说明的情况：}
  \subsection[申请人同年申请不同类型的国家自然科学基金项目情况（列明同年申请的其他项目的项目类型、项目名称信息，并说明与本项目之间的区别与联系；已收到自然科学基金委不予受理或不予资助决定的，无需列出）。]{}
  \mnsfcExcludeLine{无。}
  \subsection[具有高级专业技术职务（职称）的申请人或者主要参与者是否存在同年申请或者参与申请国家自然科学基金项目的单位不一致的情况；如存在上述情况，列明所涉及人员的姓名，申请或参与申请的其他项目的项目类型、项目名称、单位名称、上述人员在该项目中是申请人还是参与者，并说明单位不一致原因。]{}
  \mnsfcExcludeLine{无。}
  \subsection[具有高级专业技术职务（职称）的申请人或者主要参与者是否存在与正在承担的国家自然科学基金项目的单位不一致的情况；如存在上述情况，列明所涉及人员的姓名，正在承担项目的批准号、项目类型、项目名称、单位名称、起止年月，并说明单位不一致原因。]{}
  \mnsfcExcludeLine{无。}
  \subsection[申请人和主要参与者同年以不同专业技术职务（职称）申请或参与申请科学基金项目的情况（应详细说明原因）。]{}
  \mnsfcExcludeLine{无。}
  \subsection[其他。]{}
  \textbf{生成式人工智能使用情况声明}

  本人在撰写本项目申请书过程中，使用生成式人工智能工具辅助进行文献检索与整理，并对部分文字进行语言润色；所有AI生成内容均已通过原始数据库核实其真实性和准确性。申请书核心内容（包括立项依据、研究目标、内容、技术路线及创新点等）均由本人独立撰写完成，未使用AI直接生成完整申请书，亦未使用未经核实的生成内容。本人对申请书全部内容承担全部责任。
  %% 字数统计（仅在 draft 模式下显示）
  \mnsfcWordCount
\end{document}
